\pagenumbering{arabic}
\thispagestyle{plain}
\begin{center}
{\fontsize{16pt}{20pt}\selectfont\bfseries 基于几何约束与滚动规划的无线电干扰源定位与清除\par}
\vspace{5pt}
\end{center}
\vspace{12pt}
\begin{center}
{\fontsize{14pt}{17pt}\selectfont\bfseries 摘\quad 要\par}
\end{center}
\vspace{4pt}

本文研究四足机器狗在 $\pm1^\circ$ 测角误差下对无线电干扰源的自主定位与清除问题。以历史观测形成的\textbf{位置可行域}为基础，依据几何条件构造候选动作，通过预计耗时比较选择下一步，并在执行后更新可行域和任务安排，形成统一的\textbf{滚动规划}框架。

\textbf{针对问题一，}每次测向的 $\pm1^\circ$ 误差范围对应一个楔形区域，本文将其表示为两个线性不等式。$n$ 次测向产生 $2n$ 个不等式，用\textbf{半平面交}求出公共区域，即凸多边形定位区域。由凸组合性质证明凸多边形内任意两点的距离不超过最远顶点对的距离，因此区域直径可通过遍历顶点对求得。进而构造三个检测点使半平面交恰好为等腰三角形（底边 $20\,\mathrm{m}$、高 $15\,\mathrm{m}$），上顶点到直径圆圆心的距离为 $15\,\mathrm{m}$，超出半径 $5\,\mathrm{m}$，证明以直径为直径的圆不一定能覆盖定位区域。

\textbf{针对问题二，}第二检测点须同时保证能收到信号和有效缩小定位区域。本文利用接收半径 $\ge1000\,\mathrm{m}$ 的下界，\textbf{结合凸组合性质}将"到区域内每个可能位置的距离都不超过 $1000\,\mathrm{m}$"简化为对有限个顶点的距离检验，由此构造可靠接收区（各顶点对应的 $1000\,\mathrm{m}$ \textbf{圆盘的交集}），并给出构造性证明：沿示向方向约 $750\,\mathrm{m}$ 处存在半径约 $250\,\mathrm{m}$ 的圆盘整体落在可靠接收区内。在可靠接收区内沿前进方向和定位区域长轴方向分别布设候选点（横向偏移量随区域尺度自适应调整），对每个候选点在若干假想源位置上模拟测向，以移动时间、后续接近距离与剩余定位代价之和评分选点。

\textbf{针对问题三，}$10$ 至 $16$ 个全向源的数量、位置和接收半径均未知。本文设计 $8$ 个检测站（原点加半径 $999\,\mathrm{m}$ 圆周上的 $7$ 个等角点），并利用余弦不等式证明场地内任一点到最近检测站的距离小于 $1000\,\mathrm{m}$，从而保证不遗漏任何源。将剩余检测站与已发现源合并为动态任务集合，每一步的任务排序建模为\textbf{旅行商问题}并用\textbf{多起点 2-opt} 求近似最短路线；每执行一项任务即根据真实反馈更新定位区域和任务集，重新规划下一步。对已发现源，在可靠接收区内枚举测点，以位置样本和\textbf{三点 Gauss 求积}对假想测角误差加权评估后续费用；包围圆半径仍过大时，在其圆心处补测一次示向度，可证明每次这样做后半径上界缩小至原来的 $1/(2\cos\bar\varepsilon)$ 倍，保证经有限步后必定满足清除条件。区域满足清除条件时，在各顶点对应圆盘交集构成的安全区域内选择清除落点。$400$ 个本地场景和 $20$ 个模拟器演练均全部清除，平均耗时 $244.0$ 和 $225.6\,\mathrm{s}/\text{源}$。

\textbf{针对问题四，}部分源为定向发射（只向发射方向两侧各 $90^\circ$ 辐射），无信号不再等价于距离超限；贴边界朝外发射的源在场地内部无法接收。本文用 $21$ 个检测站（原点 $+$ 内环 $8$ $+$ 外环 $12$），要求每个场地点被其 $1000\,\mathrm{m}$ 内各站的\textbf{凸包}包含，使得过该点的任意半平面内都至少有一个站，从而无论发射方向如何都能收到信号。定位阶段沿示向方向布设成对检测点，\textbf{证明若两点均无信号则源沿该方向的距离不超过某一上界}，以\textbf{半平面}截断定位区域；进一步，当某无信号点到区域内每一点的距离均小于估计的接收半径下界时，可确定它在发射范围之外，由此同时约束源的位置和发射方向。沿用第三问的滚动规划，用\textbf{蒙特卡洛 rollout}比较候选检测动作的后续耗时。途中是否停下补测由 $160$ 棵\textbf{梯度提升回归树}（以 $57$ 个状态特征拟合实际耗时节省，训练集来自 $2380$ 个模拟场景）预测决定。$1140$ 个场景均全部清除，其中 $1000$ 个随机场景的平均耗时为 $449.9\,\mathrm{s}/\text{源}$。

\vspace{12pt}
\noindent\textbf{关键词：半平面交；主动测点选择；旅行商问题；滚动规划；蒙特卡洛 rollout；梯度提升回归树。}
